\section{Build vs. Buy: el límite de la diferenciación y el control plane global}

Las empresas de plataforma caen con facilidad en el error de pensar «como la infraestructura es importante, hay que desarrollarla toda internamente». Rauch sostiene lo contrario: las regiones, la red troncal, el almacenamiento de objetos y la escala de operación física de AWS son difíciles de replicar, por lo que Vercel debería acumular su diferenciación en la application-aware layer. Build vs. Buy no es una prueba de capacidad técnica, sino un problema de asignación de capital, tiempo y efectos de aprendizaje.

\subsection{Desarrollar internamente donde el insight se acumula con interés compuesto}

La sección anterior ya explicó la brecha entre las hyperscaler primitives y la application intent como un «problema de la capa de compilación»; esta sección responde además la pregunta sobre los límites organizacionales: qué capas vale la pena dominar por cuenta propia y cuáles conviene apoyar en proveedores maduros. Al leer la figura, primero compara la commodity scale con el application-aware insight y después evalúa si un desarrollo interno producirá de forma sostenida conocimiento de producto, capacidades del plano de control y ventajas de economía unitaria.

\lecturefigure{06-build-vs-buy.png}{Reutilizar la commodity scale y construir capacidades especializadas en la application-aware layer.}{Subtítulos locales 07:40--10:30; redibujo conceptual.}

\begin{knowledgebox}{Lectura de la figura: los límites de la analogía con Apple silicon}
En clase se usó el caso de Apple, que terminó diseñando sus propios chips, como analogía de la vertical integration: solo cuando una capa se vuelve central para la experiencia del producto y para la economía de largo plazo, y la organización cuenta con la escala y la capacidad suficientes, puede ser razonable integrarse hacia abajo. De la analogía no se deduce directamente que Vercel deba construir centros de datos; solo muestra que la build/buy boundary se desplaza conforme cambia la estrategia.
\end{knowledgebox}

\subsection{La metadata global es la verdadera diferenciación de la plataforma}

Vercel se parece más a un CDN/control plane application-aware que a EC2. Cada deployment es un objeto inmutable; el domain binding es un puntero mutable; el edge global necesita obtener con rapidez el pointer y la routing metadata más recientes. Una reversión no requiere sobrescribir el despliegue anterior: basta con regresar el puntero del domain a una versión conocida.

\lecturefigure{07-global-metadata.png}{Los despliegues inmutables y el domain pointer mutable separan la artifact identity de las decisiones sobre el tráfico.}{Subtítulos locales 10:00--16:40, 33:30--34:05; redibujo conceptual.}

\begin{importantbox}{Control plane y data plane}
El \term{control plane} decide deployment, domain, routing, policy y rollout; el \term{data plane} es el que realmente atiende las solicitudes de los usuarios, ejecuta funciones y devuelve assets. Las actualizaciones del plano de control deben propagarse con rapidez, pero un punto único de falla no debe interrumpir las rutas ya conocidas; el plano de datos debe mantener un comportamiento seguro cuando la metadata esté temporalmente desactualizada.
\end{importantbox}

\subsection{Ventajas y desventajas de una opinionated platform}

Las convenciones del framework proporcionan guard rails que permiten a la plataforma inferir routing, cache y runtime. El costo es que se restringe parte de la libertad, y las workloads atípicas pueden necesitar una escape hatch. Una opinionated platform sana debe hacer que la ruta habitual sea mínima y, al mismo tiempo, conservar el artifact contract, los protocolos estándar y los overrides, para evitar que el ecosistema quede bloqueado.

\lecturefigure{08-opinionated-platform.png}{Las convenciones del framework crean para la plataforma una superficie estructurada que puede optimizar.}{Materiales oficiales de Vercel FDI; subtítulos locales 13:30--17:15.}

\begin{warningbox}{Next.js no es prueba suficiente de «encerrar a los usuarios en Vercel»}
Entre el framework y la plataforma sí existe una retroalimentación positiva y un interés comercial, pero para juzgar el lock-in hay que examinar el código abierto, las salidas estándar, la vía de self-host, la portabilidad de los datos y el costo de sustitución. Estas notas analizan el mechanism y no sustituyen el análisis con etiquetas emocionales del tipo «Trojan horse».
\end{warningbox}

\teachervoice{En clase no se describió la relación entre Vercel y AWS como una competencia de suma cero. AWS ofrece hyperscale primitives confiables, y Vercel conecta con esas primitives, mediante un modelo de nivel más alto, las aplicaciones que todavía no han migrado a la nube. Las plataformas compiten por su parte del valor, pero también pueden ampliar mutuamente el mercado.}

\subsection{Resumen de la sección}

Build vs. Buy debe girar en torno a la diferenciación y al aprendizaje que se acumula con interés compuesto; las capacidades especializadas de Vercel se concentran en el framework compiler, la global metadata, el routing y el developer loop, no en repetir la escala física de los hyperscalers.
